\subsection{Warum eine Orientierungseinheit?}

Jeder studienwütige Mensch hat früher voller Enthusiasmus am ersten Tag seines
Daseins als Student bzw. Studentin das Universitätsgebäude betreten und wusste
nicht: 

\begin{itemize}\itemsep 0pt
    \item Wo muss ich hin? Was soll ich tun? 
    \item Wie sieht mein Studium aus?
    \item Welche Vorlesungen sind zu besuchen?
    \item Wo kann ich mich dafür anmelden?
    \item Wie bekomme ich einen Schein?
    \item Wer sind meine Kommilitonen?
    \item Wo gibt es Essen?
    \item Wer kann mir bei Fragen und Problemen helfen?
\end{itemize}

Durch diese Orientierungslosigkeit ergab sich bei vielen Unmut oder
Frustration. Als Ergebnis haben manche das Studium vorzeitig abgebrochen oder
einige Semester an Zeit verloren, bevor sie alle Informationen beisammen
hatten, die für ihr Studium wichtig waren.

Um solch eine Situation zu verhindern, findet die Orientierungseinheit statt,
in der wir dir diese und noch viele andere Fragen beantworten oder zu
beantworten helfen.  Außerdem wollen wir 

\begin{itemize}\itemsep 0pt
    \item dir ein paar praktische Tipps zu deinem Studium geben,
    \item mit dir exemplarisch eine Studienwoche durchgehen,
    \item dir die Möglichkeit geben, deine Kommi\-litonen kennenzulernen, 
    \item dir zeigen, was sich alles hinter dem Begriff \glqq Mathematik\grqq\
          verbergen kann,
    \item dir alles mit einer Portion Humor und Witz servieren, sodass es allen
          Spaß macht.
\end{itemize}

Dabei solltest du wissen, dass die Orientierungseinheit von Studenten für
Studenten gemacht wird.  Alle guten Ratschläge und Tipps sind also kampferprobt
und gut gemeint.  Andererseits sind aber auch Angestellte der Hochschule dabei,
die darüber wachen, dass keiner Blödsinn erzählt. Des Weiteren ist jeder von
uns (Tutor oder Veranstalter) dabei, weil er eine solche Starthilfe für wichtig
hält.

Deine Teilnahme an der Orientierungseinheit ist freiwillig, das heißt du kannst
alternativ auch die gesamte Zeit zuhause bleiben.

\columnbreak

Dann fehlen dir aber wichtige
Informationen, die dir den Einstieg in das Studium erleichtern, du lernst keine
netten, neuen Leute kennen und verpasst eine Menge Spaß. Außerdem musst du im
Laufe deines Studiums sowieso an einer Studienfachberatung teilgenommen haben,
sonst erhältst du keinen Abschluss -- warum also nicht an dieser?!

\end{multicols}
\begin{center}
\vfill
\includegraphics[width=195mm]{images/tutors.jpg}

\bigskip

\begin{minipage}{195mm}
Spaß bei der ganzen Orientierungseinheit wünschen dir

\hfill {\it die Tutorinnen, Tutoren und Veranstalter der Orientierungseinheit.}
\end{minipage}
\end{center}
\begin{multicols}{2}
